\documentclass[11pt,letterpaper]{article}
\usepackage[T1]{fontenc}
\usepackage{lmodern,microtype}
\usepackage[margin=.7in,headheight=15pt]{geometry}
\usepackage{amsmath,amssymb,mathtools}
\usepackage{xcolor,booktabs,array,enumitem,fancyhdr,hyperref}
\definecolor{ink}{HTML}{222834}
\definecolor{indigo}{HTML}{25355D}
\definecolor{muted}{HTML}{505A6B}
\definecolor{keyink}{HTML}{205D52}
\color{ink}
\hypersetup{colorlinks=true,linkcolor=indigo,urlcolor=indigo,
  pdftitle={Probability Midterm Review Quiz: Through Continuous Distributions},
  pdfauthor={Igor Rivin}}
\setlength{\parindent}{0pt}
\setlength{\parskip}{4pt}
\setlength{\abovedisplayskip}{6pt}
\setlength{\belowdisplayskip}{6pt}
\clubpenalty=10000
\widowpenalty=10000
\setlist[enumerate]{label=(\alph*),leftmargin=1.7em,
  itemsep=4pt,topsep=5pt,parsep=2pt}
\newcommand{\Prob}{\mathbb P}
\newcommand{\E}{\mathbb E}
\DeclareMathOperator{\Var}{Var}
\newcommand{\edition}{\ifanswers Instructor key\else Student quiz\fi}
\newcommand{\partpts}[1]{\unskip\nobreak\hspace{.3em}%
  \mbox{\small\color{muted}(#1 points)}}
\newcommand{\credit}[1]{\par{\footnotesize\color{muted}\textit{Credit:} #1}}
\newcommand{\problem}[4]{\par\addvspace{8pt}%
  {\large\bfseries\color{indigo}#1.\ #2}\hfill
  {\small\color{muted}\textbf{#3 points} $\cdot$ about #4 minutes}\par\nobreak}
\NewDocumentEnvironment{solution}{O{.8in}+b}{%
  \ifanswers
    \par\begingroup\small\color{keyink}\textbf{Solution.} #2\par\endgroup
  \else
    \par\nobreak\vspace*{#1}
  \fi
}{}
\pagestyle{fancy}
\fancyhf{}
\fancyhead[L]{\small\color{muted}Probability \textbar\ Midterm review}
\fancyhead[R]{\small\color{muted}Chance, Data, and Decisions}
\fancyfoot[L]{\footnotesize\color{muted}\edition}
\fancyfoot[R]{\footnotesize\color{muted}Page \thepage\ of \pageref*{quiz-last-page}}
\renewcommand{\headrulewidth}{.3pt}
\begin{document}
{\LARGE\bfseries\color{indigo}Midterm Review Quiz}\par
{\small\color{muted}50 minutes \quad $\bullet$ \quad 50 points
  \quad $\bullet$ \quad Through continuous distributions}\par
\ifanswers
\textbf{Grading note.} Award method credit after an earlier error;
avoid repeated arithmetic deductions. Suggested times total 45 minutes,
leaving 5 minutes to check work.
\else
\smallskip
\textbf{Name:}\ \rule{.49\linewidth}{.4pt}\hfill
\textbf{Date:}\ \rule{.20\linewidth}{.4pt}\par
\fi
\textbf{Instructions.} Answer all five problems and show the essential
reasoning. No calculator is needed. Exact expressions, including factorials,
exponentials, and values of the standard normal CDF $\Phi$, are acceptable.
You may quote standard distribution and moment formulas. Geometric variables
count trials \emph{including the first success}; exponential parameters
are \emph{rates}. Reserve the final 5 minutes for checking.

\problem{1}{A random permutation}{10}{9}
Four travelers' luggage tags are shuffled uniformly among all $4!$
assignments, one tag per traveler. Let $I_i$ indicate that traveler $i$
gets the correct tag, and let $X=I_1+I_2+I_3+I_4$.
\begin{enumerate}
\item Find $\Prob(X=0)$ using inclusion--exclusion. Show your count or sum.
\partpts{4}
\begin{solution}[.9in]
Let $A_i=\{I_i=1\}$. Fixing a specified set of $k$ travelers leaves
$(4-k)!$ assignments. Hence
\[
\Prob(X=0)=\frac{4!-\binom41 3!+\binom42 2!-\binom43 1!+\binom44 0!}{4!}
=\frac{24-24+12-4+1}{24}=\frac38.
\]
\credit{Uniform denominator: 1; intersection counts: 1; alternating sum: 1;
evaluation: 1.}
\end{solution}
\item Find $\E[X]$ by using the indicators. Does this calculation
require the indicators to be independent? \partpts{2}
\begin{solution}[.65in]
Each traveler has probability $1/4$ of receiving the right tag, so
$\E[X]=\sum_{i=1}^4\E[I_i]=4(1/4)=1$. Linearity needs no independence.
\credit{Indicator calculation: 1; no independence requirement: 1.}
\end{solution}
\item For $i\ne j$, find $\Prob(I_i=I_j=1)$ and decide whether $I_i,I_j$
are independent. Then compute $\Var(X)$, keeping the cross terms.
\partpts{4}
\begin{solution}[1.05in]
Fixing two tags leaves $2!$ permutations, so
$\E[I_iI_j]=2!/4!=1/12$, different from $(1/4)^2=1/16$.
Thus the indicators are dependent. Using $I_i^2=I_i$ gives
\[
\E[X^2]=4\cdot\frac14+2\binom42\frac1{12}=2,
\qquad \Var(X)=2-1^2=1.
\]
\credit{Pair probability: 1; independence comparison: 1;
second moment with cross terms: 1; variance: 1.}
\end{solution}
\end{enumerate}

\newpage
\problem{2}{One hidden choice, two observations}{10}{8}
A machine is selected once: machine $A$ with probability $1/3$ and machine
$B$ with probability $2/3$. The selected machine produces two items.
Conditional on the selected machine, defect indicators for the items are
independent. Each item from $A$ is defective with probability $1/4$;
each from $B$ is defective with probability $1/8$.
Let $D_1,D_2$ be the events that the first and second items are defective.
\begin{enumerate}
\item Find $\Prob(D_1)$. \partpts{3}
\begin{solution}[1.25in]
By total probability,
\[
\Prob(D_1)=\frac13\cdot\frac14+\frac23\cdot\frac18
=\frac1{12}+\frac1{12}=\frac16.
\]
\credit{Two weighted branches: 2; total: 1.}
\end{solution}
\item Given that the first item is defective, find the probability that
machine $A$ was selected. \partpts{3}
\begin{solution}[1.4in]
Bayes' rule gives
\[
\Prob(A\mid D_1)=\frac{\Prob(D_1\mid A)\Prob(A)}{\Prob(D_1)}
=\frac{(1/4)(1/3)}{1/6}=\frac12.
\]
\credit{Correct Bayes numerator and denominator: 2; evaluation: 1.}
\end{solution}
\item Find $\Prob(D_2\mid D_1)$. Are $D_1,D_2$ independent without
conditioning on the machine? Justify your answer numerically. \partpts{4}
\begin{solution}[1.8in]
Conditional independence within each machine and the updated weights give
\[
\Prob(D_2\mid D_1)=\frac12\cdot\frac14+\frac12\cdot\frac18
=\frac3{16}.
\]
Unconditionally $\Prob(D_2)=1/6$, by the same calculation as (a).
Since $3/16\ne1/6$, the events are dependent.
Equivalently $\Prob(D_1\cap D_2)=(1/3)(1/4)^2+(2/3)(1/8)^2=1/32$,
which differs from $1/36$.
\credit{Conditional mixture or joint-event method: 2; value: 1;
numerical independence comparison: 1.}
\end{solution}
\end{enumerate}

\newpage
\problem{3}{Counts, waits, and a complete collection}{10}{9}
Each pack contains one of three sticker types, independently and uniformly.
Packs may duplicate previously seen types.
\begin{enumerate}
\item Let $N$ count type-1 stickers in the first four packs. State its
distribution and parameters, and find $\Prob(N=2)$ and $\E[N]$.
\partpts{3}
\begin{solution}[1.25in]
$N\sim\operatorname{Bin}(4,1/3)$, with
\[
\Prob(N=2)=\binom42(1/3)^2(2/3)^2=\frac8{27},
\qquad \E[N]=\frac43.
\]
\credit{Family and parameters: 1; binomial probability: 1; mean: 1.}
\end{solution}
\item Exactly two distinct types have been collected. Let $G$ be the
number of \emph{additional} packs through the first pack containing the
missing type. State its distribution and parameter, and find $\E[G]$
and $\Prob(G>2)$. \partpts{3}
\begin{solution}[1.35in]
Each additional pack contains the missing type with probability $1/3$,
so $G\sim\operatorname{Geom}(1/3)$, on $1,2,\ldots$.
Thus $\E[G]=3$ and $\Prob(G>2)=(2/3)^2=4/9$.
\credit{Geometric model with parameter: 1; mean: 1; tail: 1.}
\end{solution}
\item Starting with no stickers, let $T$ be the number of packs through
the first complete collection. Split $T$ into the waits for the first,
second, and third distinct types, and compute $\E[T]$. Explain why
$T$ itself is not one geometric random variable. \partpts{4}
\begin{solution}[1.8in]
Let $G_k$ be the wait for a new type after exactly $k$ types have been seen,
for $k=0,1,2$. The success probabilities are $1,2/3,1/3$ and the
means are $1,3/2,3$. Since $T=G_0+G_1+G_2$, linearity gives
\[
\E[T]=1+\frac32+3=\frac{11}{2}.
\]
The chance of obtaining a new type changes from stage to stage; also
$T\ge3$, whereas a geometric variable on $1,2,\ldots$ has positive
probability of being $1$. Stage independence is not needed for this mean.
\credit{Stage probabilities: 1; means: 1; sum: 1;
reason the full collection time is not geometric: 1.}
\end{solution}
\end{enumerate}

\newpage
\problem{4}{Choose the law and use its parameters}{10}{9}
\begin{enumerate}
\item Calls arrive with Poisson counts: an interval of length $t$ hours
contains $\operatorname{Poisson}(3t)$ calls, and disjoint intervals have
independent counts. Let $M$ count calls in the next 20 minutes.
State its distribution and parameter, and find $\Prob(M\ge2)$.
\partpts{3}
\begin{solution}[1in]
Twenty minutes is $1/3$ hour, so $M\sim\operatorname{Poisson}(1)$.
Then $\Prob(M\ge2)=1-e^{-1}-e^{-1}=1-2e^{-1}$.
\credit{Family and converted parameter: 1; complement: 1; value: 1.}
\end{solution}
\item For the same call model, let $W$ be the time in hours until the
next call. Use the probability of zero calls to find $\Prob(W>t)$ for
$t\ge0$. Identify the distribution of $W$ and find
$\Prob(W>2\mid W>1)$. \partpts{3}
\begin{solution}[1.1in]
$\Prob(W>t)=\Prob(\operatorname{Poisson}(3t)=0)=e^{-3t}$.
Thus $W\sim\operatorname{Exp}(\text{rate }3)$, and
$\Prob(W>2\mid W>1)=e^{-6}/e^{-3}=e^{-3}$.
\credit{Zero-count argument and survival function: 1;
family and rate: 1; conditional probability: 1.}
\end{solution}
\item A separate task's duration $Y$ is normal with mean $10$ and
\emph{variance} $4$. Express $\Prob(8<Y<13)$ using $\Phi$.
\partpts{2}
\begin{solution}[.9in]
The standard deviation is $2$, so
$\Prob(8<Y<13)=\Phi(3/2)-\Phi(-1)$.
\credit{Standardization with standard deviation $2$: 1;
CDF difference with correct endpoints: 1.}
\end{solution}
\item Let $U$ be uniform on $[0,4]$. Find $\Prob(U>3\mid U>2)$ and
compare it with $\Prob(U>1)$. Does $U$ have the memoryless property?
\partpts{2}
\begin{solution}[.85in]
$\Prob(U>3\mid U>2)=(1/4)/(2/4)=1/2$, whereas
$\Prob(U>1)=3/4$. The unequal probabilities disprove memorylessness.
\credit{Conditional probability: 1; comparison and conclusion: 1.}
\end{solution}
\end{enumerate}

\newpage
\problem{5}{From a density to a bound}{10}{10}
A continuous random variable $X$ has density
\[
f(x)=\begin{cases}cx,&0<x<2,\\0,&\text{otherwise}.
\end{cases}
\]
\begin{enumerate}
\item Determine $c$. \partpts{2}
\begin{solution}[.9in]
Normalization requires $1=\int_0^2 cx\,dx=2c$, so $c=1/2$.
The resulting density is nonnegative.
\credit{Normalization integral: 1; value: 1.}
\end{solution}
\item Find the CDF $F(x)$ for every real $x$, and use it to find
$\Prob(1<X<2)$. \partpts{2}
\begin{solution}[1.15in]
\[
F(x)=\begin{cases}0,&x\le0,\\x^2/4,&0<x<2,\\1,&x\ge2.
\end{cases}
\qquad \Prob(1<X<2)=1-F(1)=\frac34.
\]
\credit{Complete piecewise CDF: 1; interval probability: 1.}
\end{solution}
\item Compute $\E[X]$, $\E[X^2]$, and $\Var(X)$. \partpts{4}
\begin{solution}[1.4in]
\[
\E[X]=\int_0^2 x\frac{x}{2}\,dx=\frac43,\qquad
\E[X^2]=\int_0^2 x^2\frac{x}{2}\,dx=2.
\]
Therefore $\Var(X)=2-(4/3)^2=2/9$.
\credit{First moment: 1; second moment: 1;
variance identity: 1; variance value: 1.}
\end{solution}
\item Use Chebyshev's inequality to give an upper bound for
$\Prob(|X-4/3|\ge2/3)$. State why the inequality applies.
\partpts{2}
\begin{solution}[1in]
The variable has finite mean and variance (it is bounded). Chebyshev gives
\[
\Prob(|X-4/3|\ge2/3)\le\frac{2/9}{(2/3)^2}=\frac12.
\]
\credit{Finite-variance hypothesis and correct inequality: 1;
evaluated bound: 1. This is an upper bound, not an asserted exact probability.}
\end{solution}
\end{enumerate}
\label{quiz-last-page}
\end{document}
